copyright Travis West 2025

Originally authored during my comps

I forget if this made it into my thesis or not...

\documentclass[main.tex]{subfiles}

\begin{document}    
	
	\newpage{}
	%\setcounter{page}{0}
	\pagenumbering{gobble}% Remove page numbers (and reset to 1)
	
	\begin{singlespace}
        \chapter{Integrating Datasets in the Dataflow Paradigm}\label{chapter:setflow}

        Several mapping tools have been proposed to facilitate the prototyping of DMIs. These mostly fit into two main categories: datasets (interpolations) and data flow (parameter mapping). Each has its own strengths and weaknesses, their choice likely impacting the final mappings done. [20 – 25 pages]

        a.Describe in detail 2 examples of tools in each category and elaborate on the possible impacts of each tool in the mappings obtained with them;

        b.Discuss how a hypothetical solution could combine both approaches and present a pseudo-code that could implement such a tool for one specific NIME context (programming environment, interface, synthesis method);

        c.Comment on whether these types of hybrid tools could be generalized to multiple music interaction contexts.

	\end{singlespace}
	
	\newpage{}
	\pagenumbering{arabic}
	
\section{Mapping Design Goals are Perception Design Goals}

Much of the existing literature is occupied with details of implementation,
comparing mappings in terms of whether they are explicitly or implicitly
defined \cite{Hunt2000, Arfib2002}, in terms of their functional mathematical
properties \cite{nort2004nime_geometric-properties}, or terms of their structure \cite{Momeni2006}.
While important contributions, these prior works offer relatively little
insight into how mappings are actually made, and provide no guidance on how to
make mapping design choices to achieve certain aims.  At best, they provide
vocabulary for discussing, comparing, and understanding mapping designs, but
these analytic tools are only a first step towards recognizing the properties
of mappings, not determining which properties are useful and in which
circumstances.

Van Nort and colleagues \cite{Nort2010, nort2014cmj} in addition to considering
the structural and functional properties of mappings, also emphasized the
importance of the perceptual qualities of mappings. In prior research by my
colleagues and I \cite{making-mappings-design-process,
making-mappings-design-criteria}, our preliminary observations suggest that the
perceptual view of mappings is most relevant to mapping designers. This
emphasis on perceptual considerations is also seen in prior studies observing
mapping designers: composers working with the Wekinator in a participatory
study by Rebecca Fiebrink reported that they appreciated the way the tool
allowed them to focus on gesture and sound \cite{Fiebrink2011}; and in a study
by Dom Brown and colleagues, most participants described their mappings in
terms of spatial and musical metaphor \cite{brown2018moco}.

Based on these results and my own practical experience designing mappings, it
is my view that shaping the perceptual influence of the mapping is the final
goal of mapping design. Consideration of algorithms, structure, and topology
are only practical insofar as they help us to understand the perceptual
influence of our mapping design choices. The ideal mapping design tool would
allow users to directly specify the desired perceptual results of the mapping,
and would reliably and faultlessly achieve their specification. This ideal is
unfortunately not achievable today, and realistically never will be. This is
mainly because "the desired perceptual results of the mapping" may not be known
in advance. Indeed, surprise may be an essential goal of the design, in which
case a priori specification of the behavior of the mapping would be especially
challenging. Nevertheless, the end goals of any mapping design necessarily
consider the perceptual effects of the mapping; everything else is instrumental
(in the sense of "instrumental goals").

\section{A Review of Mapping Authoring Tools}

Since the perceptual effects of the mapping cannot be directly specified or
manipulated, the mapping must be designed in terms of more immediate practical
objects. I conducted a systematic review of mapping design tools presented at
NIME and related venues, considering the ways in which mappings are designed,
as manifested in these design tools. Based on my review, mappings are most
often designed in terms of dataflow and/or datasets; that is to say, the
perceptual behavior of the mapping is influence indirectly by manipulating the
flow of data and/or the curation of datasets. Design tools can be characterized
in terms of their affinity for dataflow and datasets, among other features.
Based on these considerations, there are clear directions in which mapping
design tools may be improved, and ways in which new tools may build upon
previous work to better empower mapping designers.

\section{Method}

An inital pool of literature was formed starting with the special issues on
mapping from the journals Organised Sound (year 2002, issue 2, volume 7) and
Computer Music Journal (year 2014, issue 3, volume 38) as well as the
bibliography of the ISIDM \cite{isidm-bibliography}. The ISIDM bibliography
does not include any papers after 2011, so I additionally included all papers
published at the NIME conference from 2012 to 2022 (excluding papers
I wrote) as a representative body of research since 2011. In addition, a few
papers were included for analysis from my personal bibliography based on my
prior familiarity with the literature.

% numer of tools: $ readings list | xargs rg --no-heading -NI '@@mapping@@design@@(sub-)?tool:' | sort | uniq | wc -l
% number of papers: $ rsa '@@mapping@@design@@(sub-)?tool:' | wc -l

The initial pool was scoured for papers presenting mapping design tools using a
combination of keyword search and manual skimming.  62 mapping tools were
identified based on 84 papers, including 49 papers presenting a design tool.  I
read each of these papers and annotated them considering the following
questions:

\begin{itemize}
    \item What are the objects that users are concerned with and manipulate?
    \item How do users interact with these objects of interest?
\end{itemize}

\section{Dataflow Mapping Design Tools}

% number of tools   $ rsa @@mapping@@design@@tool: | rs @@orientation@@dataflow | readings col-tags | rg @@tool: | sort | uniq | wc -l


It is no surprise that dataflow-oriented design tools feature prominently in my
review, including 23 of 49 tools surveyed. As described in an earlier chapter,
the earliest digital music systems borrowed the flow of signals as a guiding
principle and conceptual framework from their predecessors in analog modular
synthesizers and tape music studios. Modular approaches to the design of music
software absolutely dominate commercial music production software, and most
music programming environments like Max/MSP, Pure Data, Supercollider, and
Chuck, all support this conceptual framework.

% number of libraries   $ readings list | xargs rg @@.*library

Precisely what to call the objects of interest in a dataflow mapping design
model differs from author to author. I will adopt the terminology that I
previously suggested in my MA thesis \cite{West2020thesis}.

In a dataflow model, \emph{devices} are the owners of \emph{endpoints}.  An
endpoint may be the \emph{source} of a \emph{signal}, or it may be the
\emph{destination} where a signal is sent to influence a \emph{parameter} of
the device.  When a signal is routed from a source to a destination, this
constitutes a \emph{connection} or an \emph{association} between the two
endpoints.

As well as endpoints, devices, and connections, dataflow mapping designers are
often also interested in controlling the \emph{range} of an endpoint or
connection, i.e. the minimum and maximum value of an endpoint. Designers often
design associations from the range of one endpoint to the range of another;
this linear transformation is often viewed as a part of the connection, as in
Webmapper \cite{Wang2019} for example. More generally, the designer may be
interested in the overall transfer function implemented by a connection. This
view of connections as providing necessary transformations to allow one endpoint
to be compatible with another is common.

\section{Objects of Interest in Dataset Mapping Design Tools}

Machine learning tools have been used in the design of mappings even before the
importance of the design of mappings began to receive emphasis in the
literature.  An early example is Lee and colleagues' addition of artificial
neural networks to an early version of Max (before MSP was added), presented in
1991 \cite{lee1991ann}.  Since then, numerous different models and algorithms
have been gainfully employed in the design of mappings, including 19 of the 49
tools surveyed.

In these tools, users main interactions with the mapping are achieved by
creating or deriving data that implicitly encode the specification of the
mapping. The data may take the form of point-wise demonstrations, common in
preset interpolators, recordings labelled with classification, often generated
by hand, and other multi-modal data recordings, such as those generated by
miming gesture along with sound sequences playing in real time. Several
projects have also proposed ways of generating datasets in real time without
user intervention, leading to systems that partially or fully automatically
learn the mapping without supervision \cite{smith2011self-supervising,
scurto2017online-clustering, delaspozas2020nime,
murray-browne2021nime_latent-mappings}.

The datasets, however they are derived, are used to train adaptive systems that
model the relationships illustrated by the dataset.  Running the models then
allows output data in one domain to be generated by inputting data from a
different domain, e.g. sound synthesis parameters can be determined from sensor
signals.

At a fundamental level, the objects of interest when using these tools are
points of data, streams of recorded or generated data, and associations linking
data (points or streams) from different domains to form multi-modal datasets.
The models themselves may be manually edited to some extent, but this is done
so that the model will more accurately learn the relationships reflected in the
dataset.

\section{Higher-order Mapping Design and Mixed Data-flow/-set Approaches}

If a first-order mapping associates gesture signals to sound parameters, then a
mapping which controls a first-order mapping may be considered a higher-order
mapping. This kind of structure is usually called a \emph{convergent} mapping,
although this term has also been interpreted to be synonymous with a
\emph{many-to-one} or \emph{many-to-many} structure, which is not necessarily
higher-order in the sense described \cite{West2020thesis}. When dataflow- and
dataset-oriented approaches are combined, this may also lead to higher-order
mappings.

In many cases, the mappings achieved with dataflow tools are affine mappings,
and can be represented by an augmented transformation matrix.  In such cases,
the mapping matrix itself can be treated as a device whose coefficients are
parameter endpoints that can be modified in real time; a higher-order dataflow
technique.  This approach is suggested by de Campo \cite{decampo2014icmc_metacontrol}
in the form of applying random offsets to matrix coefficients.

On the other hand multiple matrices can be pre-set and interpolated between to
dynamically alter the mapping.  In this approach, the mappings themselves are
treated as data.  They can be associated with data points in a different
domain, such as 2D spatial coordinates, and an algorithm such as a preset
interpolator can be used to navigate the space of mappings by moving through
the space of 2D coordinates.  This kind of approach was found in 3 of the tools
surveyed \cite{bevilacqua2005nime_mnm, Brandtsegg2011, decampo2014icmc_metacontrol}.  Alternatively,
mappings themselves can be subjected to evolutionary techniques and subjective
breeding, as demonstrated by Mandelis and Husbands in 2004
\cite{mandelis2004grow-it}.

Matrices are also interesting as there are efficient algorithms for learning a
matrix transformation based on a dataset, but the matrix itself is also
possible to understand for designers.  As such, matrices are perhaps unique in
their ability to be directly tweaked and specified by the designer or trained
from data.  Matrices are also conceptually simple and relatively easy to
implement, which makes them particularly attractive as mapping primitives.
This approach was found in 1 of the tools surveyed \cite{bevilacqua2005nime_mnm}.

Dataset-oriented interactions can be integrated into a dataflow framework by
allowing trained models to exist as devices in the dataflow graph.  This
approach was identified in 3 of the tools surveyed \cite{bevilacqua2005nime_mnm,
gillian2011sec, Bullock2015}, but arguably can be found in more than those
three, since many dataset-oriented tools have been implemented as modules for
dataflow-oriented languages, notably Max/MSP and Pure Data, e.g.
\cite{francoise2014nime_xmm, kiefer2014esn}.  This potential for integration is
promising, but there may be significant room for usability improvements, to be
discussed below.

Many algorithms used in dataset mapping design have free parameters, especially
in preset interpolation algorithms.  In a dataflow-oriented approach, these
parameters could be treated as endpoints and signals could be connected to
influence them in real time.  In some cases, it may be practical to run a
dataset-oriented model in real-time without training it in advance.  In such
cases, the dataset itself can be treated as a collection of endpoint
parameters.  Depending on the size of the dataset, it might be practical to
individually route signals to each point in the dataset, controlling them in
real time.  This approach was not emphasized in any of the tools surveyed.

It is likely that any given dataset may be too large to treat in a
dataflow-oriented approach.  However, datasets may be treated as parameter
vectors in a higher-order dataset-oriented approach. For example, in their 2017
NIME paper, Visi and colleagues \cite{visi2017nime} derive a topological
representation of datasets recorded from multiple participants and use a common
preset interpolation technique to interpolate between datasets, generating new
intermediary datasets that are then used to train an adaptive model to follow
gestures.

\section{Wherein Lies the Design?}\label{section:endpoints-and-connections}

In dataflow-oriented mappings, knowledge about the behavior and affordances of
devices is encoded in their endpoint representations.  The user specifies the
mapping by making connections between endpoints.  The design is encoded in the
connections, and the knowledge implicit in the endpoints.  When the endpoints
encode a significant amount of knowledge, the connections may be relatively
simple; for example, there is no need for a complicated mapping between a
breath pressure sensor signal and a physical model of a flute with an air
pressure parameter.  In case the endpoints encode very little information, the
mapping may need more complexity, such as multiple processing layers, in order
to extract, condition, and combine signals so that a meaningful set of
endpoints is derived.  The knowledge embedded in a dataflow-oriented mapping
design is therefore strongly associated with the endpoints in the mapping.
This is reflected by the way that endpoints can often be named and described
in a dataflow network. That's the filter cutoff control, that's the shake gesture,
and so on.

Dataset mapping design tools can be viewed as being at least somewhat
dataflow-oriented, as once the dataset is used to train a model, the model is
naturally treated as a device through which the user will route signals. But
unlike typical devices in a dataflow approach, a mapping model trained on data
does not have its own set of meaningful endpoints. Its parameters are a
reflection of the input domain, and its signals are a reflection of the output
domain. As such, the endpoints of a model trained on data do not encode any
knowledge on their own. The trained model could even be viewed as having no
endpoints; it is more akin to a single complex connection between the (often
high-dimensional) domain and range represented by the signals coming into the
model and the parameters it influences. Instead, \emph{the dataset used to
train the model} contains knowledge about the design. Presets, gesture
recordings, and sound parameter automations have inherent meaning, especially
when they are carefully designed by the user. Linking them together to form
multi-modal datasets is analogous to making a connection between endpoints in a
dataflow model. The data in a dataset-oriented model can thus be viewed as
implicit representations of meaningful connections and endpoints that the
learning models in these tools are tasked with implementing through their
internal machinations.

If this is the case, we should expect that the datasets in dataset-oriented
techniques would be associated with meaningful concepts, which is indeed what
we observe.  In all of the dataset-oriented tools surveyed, the datasets of
interest were related to gestures and/or sound parameters that were known to be
of interest in advance. The fact that datasets and endpoints encode specific
meanings is also shown by their naming conventions. Like endpoints in a dataflow
network, datasets can often be labelled in terms of what they are meant
to represent, such as the names of presets or gesture demonstrations.

In preset interpolators, the fact that the dataset is used to derive more
meaningful endpoints from raw signals is explicit in the structure of the
algorithm. The presets in control space represent meaningful states of the
control space signals. These algorithms employ various geometric models of
closeness to these control points to generate a vector of weight signals that
are then used to mix between the presets in destination space. The vector of
weights is literally the realization of a model of meaningful endpoints
represented by the control space presets.

An exception to rule that dataset-oriented models hide their
meaningful endpoints internally is found with uni-modal tools such as gaussian
mixture models \cite{francoise2014nime_xmm}. In these tools, the model is trained to
derive a set of meaningful parameters, such as gesture descriptors, from a set or raw inputs,
and so is another case where the learning model is explicitly trained to represent
meaningful endpoints.

Therefore, in both dataflow- and dataset-oriented mapping design approaches,
endpoints (i.e. signals and parameters) may be considered as primary vessels of
knowledge and intent in the mapping design, whether they are empirically
represented or implicitly encoded in datasets and represented only in the
internals of a learning model.  Connections between endpoints,
whether explicit connections or multi-modal associations between datasets, are also an
essential component of the mapping design, since they specify how the potential
behaviors and affordances encoded in the endpoints should be related.  

However, higher-order mapping approaches blur the lines between endpoints and
connections.  For example, a transformation matrix may be thought of as a connection from
$N$ signals to $M$ parameters. It may also be viewed as a collection of
data modelling something that can be learned by a training model, i.e.
as an implicit representation of endpoints and their connections.  The same
matrix could also be viewed as an endpoint unto itself, a single $N*M$ dimensional
point to be modulated by some other process. The ambiguity here arises due
to the fact that mappings themselves can be viewed as having endpoints,
or as being defined by datasets that implicitly model those endpoints.

From these considerations, it appears that if datasets are viewed as implicit
models of endpoints and connections, then endpoints and connections are
conceptually sufficient objects of interest that can fully model mapping
designs. Notably, devices seem to play a similar role to datasets; both devices
and datasets group together endpoints, and implement hidden layers of
processing to derive a set of outputs from their inputs. Devices, like
datasets, hide internal connections that associate inputs to outputs. Unlike
datasets, the endpoints of devices are generally semantically meaningful unto
themselves, whereas datasets tend to hide some or all of their meaningful
endpoints internally, and connect directly to externally meaningful endpoints.
Datasets and devices can thus be viewed as conceptually related, or even
instances of the same conceptual object. They group together endpoints, and
implicitly connect them. They are black boxes whose internal processes are
inscrutable, or at a lower level of detail than is called for by the designer.

\section{Models and Cross-Domain Mappings}

Returning to the importance of the perceptual view of mappings, and considering
the conceptualization of a mapping as a network of endpoints and connections,
whether explicitly grouped together by devices or implicitly by datasets,
I would like to emphasize that there are two broadly different kinds of reasons
for mapping design, regardless of the tools used to implement the design.  These
two motivations can be used to broadly segregate the mappings in a DMI into
two groups.

Before addressing these reasons for mapping, consider the mappings that appear
in and around a DMI. The concept of a mapping is generic and flexible. It is
not unreasonable to broadly imagine the concept, and consider the mappings e.g.
from the performer's intention to their physical movement, from the physical
movement to the actuation of sensors, from the actuation of sensors to the
transduction of electrical signals, from the electrical signals to digital, and
so on until digital signals are converted back to electrical, electrical to
acoustic, acoustic to neurological, and from neurological signals another
mapping to meaning in the mind of the listener. The connections also extend
back from the audience to the performer, via visible light and other signals,
which are mapped in the performer's brain to an interpretation of the
audience's affective states, which may in turn influence the performer's
intention. However, all of these mappings just described are largely outside of
the purview and influence of the mapping designer. We could term them inherent
mappings or external mappings. They are largely outside the scope of the
designer's considerations, or at least outside the scope of their practical
implementations.

Within the designer's scope, there are also numerous mappings, in the broadest
sense. Signals from the sensors must be conditioned, calibrated, fused, and
analysed, perhaps producing gestural signals that can be further conditioned,
calibrated, fused, and analysed to produce other signals of interest. Rather
than being mapped directly to synthesis parameters, there might be on the
synthesizer's side further layers of conditioning, calibration, fusion, and
analysis used to produce higher-level sound parameters. In the middle, a
simpler layer of mappings might bridge the gap from high-level gesture signals
to high-level synthesis parameters, crossing from the domain of movement to
that of sound. This is the layered structure of mappings described by Hunt et
al \cite{Hunt2002a}, proposed as a way of simplifying the design process, and
potentially providing for portable mappings that could generalise to many
instruments.

The mappings giving rise to intermediate layers, in Hunt et al's discussion, are notably
different in character from the innermost layer of mappings. Let us term the latter as
\emph{cross-domain mappings}, and the layers of conditioning, calibration, fusion, analysis, etc.
not as mappings at all, but as \emph{models} of specific meaningful signals and behaviors.
In the conceptual framework of meaningful endpoints with connection between them,
models are characterized by giving rise to new meaningful endpoints, whereas cross-domain
mappings merely associate existing endpoints.

Revisiting the discussion of evaluation, ecologies, proliferation, comparison, and developing design
knowledge, it is also apparent that models exhibit better generality, relating only to the
specificities of their inputs or outputs, whereas cross-domain mappings relate to both.
Models are thus more likely to be useful in the context of multiple musician+instrument
ecologies. For example, gesture models that derive new signals from MIMU data may be useful
to any instrument with a MIMU, whereas cross-domain mappings from those gestures to synthesis
parameters are only potentially useful if an instrument uses a MIMU \emph{and} the same synthesizer,
or one exposing comparable high-level synthesis parameters.

On the other hand, recognizing the diversity and complexity of the environment, I do not advocate
that high-level signals produced by models should \emph{replace} or otherwise reduce access to
the least-conditioned raw data, nor do I wish to imply that the goals previous assumed should
be universally adopted and reflected in the design of mappings.

As described by Thorn and Sha \cite{thorn2019moco_instruments-of-articulation},
"signal processing affords continuous revision of instrument 
morphology without requiring termination in a conceptual schema 
demanded under the scientific desideratum of universal validity." They describe their approach
to the design and construction of "instruments of articulation" emphasizing a closed loop
of sensorimotor engagement with a system leading to changes in the code that influence
sensorimotor experience, a nomadic science where the development of the system responds to the
conditions the system creates, through pure signal processing that doesn't necessarily
encode a priori categories or semantic interpretations, but rather responds interactively to
events, conditions those events, and produces results from which meaning emerges rather than
enforcing assumed meaning through its design. Thorn and Sha's discussion constrasts
vividly with my own emphasis and aspirations toward generality, universal usefulness, and
interpretable networks marked with semantically laden endpoints. Let it be understood that
these approaches are not mutually exclusive. In particular, so long as the inner mechanisms
of an instrument system are exposed and accessible, Thorn and Sha's process of canalization
remains possible, and within the conceptual framework I have constructed, there is no
reason why the meaning inherent in an endpoint must necessarily submit itself to ready
interpretation, categorization, or universal fixed understanding. Endpoints can also
be semantically shallow, in the sense described by Thorn and Sha.

\begin{displayquote}
Through trial and 
error, salient features are designated within a code that constructs 
those features. Choices are made about what needs to be detected, 
this band of energy or that type of attack. The media itself is 
constructed by the actual mathematical choices in the analysis of 
the feature vectors. ... In building instruments that are radically open 
for players to rupture pre-conceived semantics through their 
improvised, alinguistic and articulating actions, we are advocating 
a view of intentionality that is semantically shallow and not based 
on the notion of intentional being connected to a goal known in 
advance.
\end{displayquote}

\section{Deep Dive: Dataflow}

This section describes the basic resources of mapping design, which tend to
emphasize the flow of data through the system. These are essentially the
techniques of computer programming. They can be leveraged to model semantically
meaningful qualities of the signals in an instrument, or to merely condition
the connection from one endpoint to another in a cross-modal mapping.

\subsection{The Modulation Matrix}

A mapping viewed as a network of connected endpoints naturally lends itself to a mathematical treatment
as a directed graph, where the endpoints correspond to vertices and the connections to edges. This can be
represented with an adjacency matrix. This representation is sufficient for encoding simple connections
between endpoints with scalar values. Multidimensional values can be accomodated by breaking them into
scalar components. Treating the network as a weighted graph gives rise to a matrix
with scalar coefficients, allowing the influence of each connection to be adjusted.  This is one way
of viewing a modulation matrix.

If there are $N$ sources and $M$ destinations, then the matrix will have dimension $N*M$. The current values
of all the sources can be treated as a column vector $s$, and the value of the signals input to all
destinations, a column vector $d$ is given as $d = As$, where $A$ is the matrix representation of the graph.
Each non-zero element $A_{nm}$ in the matrix represents a connection from the $n$th source to the $m$th destination.
The $n$th row in the matrix can be viewed as showing the influence of the $n$th signal on every one of the parameters.
The $m$th column gives the influence of every signal on the $m$th parameter.

Augmenting the graph with a constant signal with a value of 1 allows signals to be offset as well as
adjusted in magnitude. An $N+1$th row is added to $A$ which gives the offset for each parameter.
This matrix is equivalent to an affine transformation from the source space to the destination space.

A modulation matrix of this sort is sufficiently powerful to represent many practical dataflow mapping designs.
The natural interpretation of the matrix as routing influence of the sources to the destinations also lends itself to
higher-order extensions. De Campo \cite{decampo2014icmc_metacontrol} suggests applying random perturbations
to the matrix coefficients, crossfading or blending between pre-defined matrices, and modifying the coefficients
by other means. Brandtsegg et al \cite{Brandtsegg2011} also highlight the usefulness of blending between
matrices, and Bevilacqua et al \cite{bevilacqua2005nime_mnm} also suggest its utility.

Bevilacqua et al \cite{bevilacqua2005nime_mnm} further highlight the possibility to derive matrices of this sort
through linear regression of a dataset, using singular value decomposition. This illustrates modulation matrices
as a particularly unique and powerful mapping tool, because they can be readily composed by hand with a
dataflow-oriented approach, and they can also be easily derived from a dataset. Their flexibility toward
high-order approaches makes them an especially promising candidate for facilitating integrated
flow/set agnostic approaches to mapping design.

\subsection{Other Dataflow Resources}

Since dataflow is rather the default assumption in mapping design, and to some extent in computer
music and music technology more broadly, the basic resources of dataflow mappings are not in general
emphasized in the literature as notable techniques, per-se. Rather, they are taken for granted as
basic ingredients. The modulation matrix, for example, is essentially a particular way of structuring
multiplication by a constant, addition of a constant offset, and mixing together of signals from multiple sources,
which may reasonably be viewed as such basic operations as to be beneath notice or discussion.

Other tools that can be taken for granted include basic nonlinear functions such as absolute value, exponentiation, and
logarithms, sigmoid curves such as the hyperbolic tangent, piecewise linear breakpoint envelopes, splines,
and other (sometimes piecewise) polynomial functions.

Some of these can naturally be integrated into the modulation matrix paradigm. For example, the multiplication
by a constant at each cell in the matrix can be replaced by a breakpoint envelope or other non-linear curve.
This kind of approach is found in commercial plugins such as Roli's Equator synth\footnote{\url{https://roli.com/products/software/equator2}} and mntra instruments' MNDALA engine \footnote{\url{https://www.mntra.io/mndala/}}. In principle, powers and products
of signals could be added as additional inputs into the modulation matrix, allowing it to model convergent
polynomial-type relationships. Similarly, non-linear functions could be applied to signals at the input, or
to the outputs of the matrix, integrating these into the matrix as additional endpoints on its edges.

Alternatively, and more commonly, these operations might be considered as primitive devices unto themselves.
This readily gives way to a general boxes-and-lines programming environment, with all of its attendant
capability and complexity.

\subsubsection{The Effects of Dataflow}

while highly useful, modulation matrices are also limited in their expressive power.
affine transformations cannot implement non-linear mappings, and while smaller
matrices may be developed by hand, reasoning about very large modulation matrices
can become challenging, and their design by hand cumbersome. Integrating non-linear
processes in the matrix paradigm can overcome the first issue, and the expense of
increasing the second.

More generally, dataflow oriented techniques are only as useful as the
designer's ability to understand and implement flows that result in the desired
behaviors. When design goals or desired behaviors are not clearly specified,
dataflow based design approaches offer the designer little help. Furthermore,
the requirement for the designer to largely specify all operations in the
mapping makes it challenging to implement mappings that exhibit surprising or
unexpected behavior. These limitations serve as strong motivation for the
development of dataflow-oriented design techniques.

\section{Deep Dive: Datasets}

This section introduces a handful of tools which generate mappings from data. In particular,
I will describe several preset interpolation algorithms, since these are a common kind of
dataset-oriented mapping technique that are widely used, easy to implement, and provide
interesting results.

The basic principle of preset interpolation is to mix between
$N$ presets $p_n$, which are simply vectors in the space of destination parameters,
producing an output vector $y$ in destination
space via a weighted sum with weights given by $w$:

\begin{displaymath}
y(w) = \sum_{n = 0}^{N - 1} w_n * p_n
\end{displaymath}

Notice that this application of the weights can also be viewed as a matrix
multiplication, where the matrix is given by concatenating the presets (viewed as columns),
e.g. $y = Pw$.

The role of the preset interpolator is simply producing and
applying a list of weights, usually as a function of a query or cursor vector
$q$ in the source space, and list $D$ of $N$ demonstrations, where each demonstration $d$
is simply a pair of vectors ${s_n, p_n}$ with $s_n$ in the source space and $p_n$ in the
destination space.

\begin{displaymath}
w(q, D) = [w_0, w_1, w_2, ..., w_{n}]^T 
\end{displaymath}

\begin{displaymath}
D = \{s_0, p_0\}, \{s_1, p_1\}, \{s_2, p_2\}, ..., \{s_n, p_n\}
\end{displaymath}

The complexity of the algorithm lies in the way in which the weight vector is produced.
This usually employs a geometric interpretation of the source vectors, where
some notion of "closeness" from the query vector to the source vectors is used
to derive a weight for each source vector, which is then applied to the associated
preset vector.

While numerous preset interpolation schemes have been proposed for musical applications,
I will cover two in detail which happen to be my favorites based on informal testing.
They are both reasonably simple to implement, amendable to higher-order mapping structures,
and between the two of them provide a reasonable sense of some of the concerns involved
in the selection of a preset interpolation algorithm, such as locality and smoothness.

\subsection{Inverse Weighted Distance}

Likely the simplest preset interpolation algorithm,
well described by Todoroff et al \cite{todoroff2009icmc_interpolation-tools},
follows a simple inverse distance rule to determine the weights.
The most basic version of the algorithm is a simplified gravitational model.
The weight $w_n$ for demonstration $n$ with source point $s_n$ is simply
the inverse of the squared distance from the query $q$ to $s_n$:

\begin{displaymath}
w_n = \frac{r}{||q - s_n||^2}
\end{displaymath}

Where $r$ accounts for the relative mass of the demonstration. A useful
extension is to allow an arbitrary positive exponent $p$:

\begin{displaymath}
w_n = \frac{r}{||q - s_n||^p}
\end{displaymath}

Todoroff and colleagues also introduce a global $d_{min}$ parameter that serves
to prevent numerical instability, and $r_{min}$ parameter that allows a region
around the demonstration to be dedicated to that demonstration.

\begin{displaymath}
w_n = \frac{r}{max(d_{min}, (||q - s_n|| - r_{min})^p)}
\end{displaymath}

In principle, the mass $r$, minimum distance $d_{min}$, exponent $p$ and distance offset $r_{min}$
can be varied on a per-demonstration basis.

Various extensions to this already quite flexible model are possible. A background
value can be added with a low constant weight to take over when the query vector
is distant from all control points \cite{todoroff2009icmc_interpolation-tools}.
With such a background vector, another useful extension is to ignore control
points that are further away than a certain minimum threshold distance, improving
the locality of the model \cite{gibson2019smc}.

Considering the functional characteristics of the algorithm \cite{nort2014cmj,
nort2004nime_geometric-properties}, the basic inverse weighted distance
interpolation is smooth (i.e. continuously differentiable), continuous, allows
demonstrations to be scattered (i.e. they needn't be aligned to a grid), and
N-to-M dimensional (i.e. maps from an arbitrary number of source dimensions to
an arbitrary number of destination dimensions), as well as very simple to
implement, with $O(n)$ complexity in the number of demonstrations, even if the
demonstrations are changed dynamically in a higher-order mapping structure. The
minimum distance and radius extensions described above introduce
discontinuities in the first derivative, but otherwise maintain the good
functional characteristics of the model. The main drawbacks of this algorithm
are that it exhibits some non-local behaviour; adjusting the location or other
parameters of one preset can influence the resulting output far from the
location of that preset, because the weight of distant presets is generally
never exactly 0. Gibson's minimum threshold extension \cite{gibson2019smc}
improves this situation somewhat. It could also be argued that the large number
of tuneable parameters may be a nuisance, in case the user does not wish to
have to carefully tune the response of the algorithm.

\subsection{Marier Spheres}

Martin Marier introduced the so called "intersecting n-spheres interpolation"
algorithm \cite{marier2012nime_nspheres} as part of his work with The Sponge
\cite{marier2010nime_the-sponge}, a deformable digital musical instrument which
Marier has been developing for over a decade. This algorithm produces an
interpolated output that is smooth, continuous, N-to-M, exactly passes through
the demonstrated destination points, and locally depends only on nearby
demonstrations. The algorithm also imposes no constraints on the placement of
source and destination presets, and with an efficient implementation the
algorithm allows for the position of presets to be dynamically varied in
real-time, allowing the possibility for higher-level mappings.

Intersecting n-spheres interpolation produces weights for a weighted sum
according to the following procedure. In summary, for each demonstration two
circles are drawn: one centered on the source vector of the demonstration (the
source circle), and one centered on the query vector (the query circle). The
radius of the query circle is given by the distance from the query point to its
nearest neighbour in source space. The radius of the source circle is also
given by the distance to its nearest neighbour, which may be either the query
point or a different source vector in a different demonstration. The weight for
each demonstration is given by the ratio of 1) the area of the source circle
for that demonstration, and 2) the area of the intersection of the source
circle with the query circle. In a practical implementation, the procedure in
full proceeds as follows.

Given a query point $q$, and a list of demonstrations consisting of $N$ points
in source space $s_n$ paired with the same number of points in destination
space $p_n$:

1. Determine the distance $d_n$ from each demonstrated source point $s_n$ to
the cursor point $q$. At the same time, search for the nearest neighbour of $q$
and record its distance from $q$ as $r_q$.

2. For every demonstrated point in source space $s_n$, calculate and
record the distance $r_n$ from that point to its nearest neighbour in source
space (not including $q$).

3. The weight $w_n$ associated with the destination vector $p_n$ of the $n$th
demonstration is given by $A_x(r_q, r_{min}, d_n) / A(r_{min})$ where $r_{min} = min(r_n, d_n)$, $A(r_{min})$ is
the area of the source circle, and $A_x$ is
the area of the intersection of the source and query circles.
Note that $A_x$ depends only on the radii of the two circles $r_q$ and $r_n$,
and the distance $d_n$ between them.

4. Calculate the interpolated output $p$ in destination space as the weighted
sum of demonstrated outputs $p_n$.

5. Normalize the weights (or equivalently the weighted sum) so that their sum
is equal to one by dividing each weight $w_n$ by the mean of all the weights
$\mu = \frac{1}{N}\sum_{n = 0}^{N - 1} w_n$

$A_x$ is given by the following function:

\[
\begin{aligned}
A_x(R, r, d) &= r^2 cos^{-1} ( \frac{d^2 + r^2 - R^2}{2dr} ) \\
             &+ R^2 cos^{-1} ( \frac{d^2 + R^2 - r^2}{2dR} ) \\
             &- \frac{\sqrt{(-d + r + R)(d + r - R)(d - r + R)(d + r + R)}}{2}
\end{aligned}
\]

And $A$ is given by the usual equation for the area of a circle:

\[
A(r) = \pi r^2
\]

Step 2 above is an instance of the "all nearest neighbours" problem. We are
given a set of points (the demonstrated points in the source space plus the
cursor point) and for each point $p$ we need to determine the nearest other
point in the space, which will be used to determine the
radius of the circle associated with $p$.

Based on a cursory and incomplete search, efficient
algorithms appear to exist which solve the all nearest neighbours problem in
$O(n\log n)$ time, or even in $O(n)$ time if certain constraints are met.
Unfortunately, comprehending (let alone implementing) the linear time
algorithms for this problem is non-trivial, so we opt to use a simpler
algorithm until such time as it becomes clear that optimization is necessary.

If we assume that the demonstrations may change at any time, then updating any
of the points in source space including the cursor point always requires the
all nearest neighbour problem to be solved anew, because the nearest neighbour of
the updated point and any point that it was previously the nearest neighbour of
all have to be updated. Because the position of the cursor point is assumed
to change whenever the interpolator is queried, this means that the nearest
neighbour problem must be solved for every query. A trivial implementation
such as in Marier's original Supercollider implementation of the
algorithm solves this problem with $O(n^2)$ time using plain old
brute force; the distance from every point to every other point is calculated,
and the nearest neighbour is found as the point with the least calculated
distance value. In Marier's practice this $O(n^2)$ time complexity hasn't been
reported as an issue, presumably because the number of data points has always
been reasonably low.

However, if we assume that the demonstrations have not changed since the last
query, then it should be unnecessary to solve the all nearest neighbours
problem.  Instead, the weights can be generated in linear time, since only the
distance from the query point to each demonstration needs to be calculated.

The algorithm, despite its name, actually doesn't involve any n-spheres.
Intead, 2D circles are used, simplifying the algorithm while producing results
that "are predictable and feel natural to the user" according to Marier. The
advantage of Marier's algorithm is that the influence of each demonstration is
constrained to its local neighbourhood. In my informal experiments with the
algorithm, the sense I have is that each demonstration has a certain opacity,
so that the query vector is not able to see past nearby demonstrations, so that
far away demonstrations have no influence on the output. This also has the
effect that when the query vector is far from any point, usually only the
nearest point has any influence. In comparison, inverse weighted distance
algorithm produces a muddy mix of all presets when the query vector is far from
everything. An unusual quirk of Marier's spheres occurs when two demonstrations
are very close in source space. This causes their source circles to be
excessively small in size, producing interesting artefacts and strongly local
behavior near those demonstrations. It is unclear whether this should be
considered a drawback or a useful feature.

\subsubsection{The Influence of Datasets}

The main advantage of dataset-oriented tools is that they allow the designer to
provide data that is known to encode something interesting, and the algorithm
then provides an implicit representation of that behavior. This enables
mappings to be designed despite poor a priori specification or understanding,
and also gives rise to interesting behaviors being assigned to swathes of
intermediate inputs whose corresponding outputs are not explicitly defined by
the designer. In user studies (e.g. \cite{Fiebrink2011}) this has been found to
permit designers to focus more readily on the perceptual outcomes of the
mapping processes and to produce more surprising and explorable mappings ,
whereas dataflow oriented approaches in contrast may require designers to focus
largely on low level details of implementation.

\section{Integrated Design of Mappings by Dataflows and Datasets}

Although there is relatively little empirical information on how the design process
is influenced by differently oriented mapping design tools, it can reasonably be
presumed that dataflow and dataset type approaches offer different efficiencies
and influence the design process is observably distinct directions. For this reason,
it may be advantageous to provide mapping design environments in which these
techniques are equally accessible for use in the design process. However, a number of
challenges arise when considering how to achieve this kind of integration.

In a dataflow-oriented approach, endpoints are generally considered as scalar values.
Even when naturally vectorial endpoints are involved, such as the acceleration of an
object as measured by an accelerometer or the color of a pixel, it is common to treat the elements of the
vector as individual endpoints, and to seperately process and utilize them. In general,
endpoints are practically asynchronous, due to the performance characteristics of the
transmission media used to communicate signal data from the embedded environment where
they are acquired to other processing nodes in the system. For example,
a button's signal may only be transmitted when the state of the button changes.
Where signals are actually acquired in a synchronous manner, it is common for information
about the exact sampling rate to be discarded, if it is ever known, and it cannot
in general be assumed that synchronous signals from different sources will share
any common timebase allowing them to be synchronized with each other.

In contrast, dataset-oriented approaches unanimously treat endpoints as regularly-sampled
points in (usually Euclidean) vector spaces. All, or a subset, of signals generated by
an input device are lumped together as "the source space", and similarly the parameters
of a device are treated as a vector in "the destination space". This is difficult efficiently to
reconcile with the more fine-grained view adopted in dataflow approaches, as well
as the typically asynchronous and event-driven nature of most implementations.
Naive approaches raise concerns about costly re-computation of values, similar to those
that arise in functional reactive programming systems; if a single element of a vector
in a dataset is updated, this may trigger recomputation of a complex data structure
internal to a dataset-driven system, such as the costly all nearest neighbours search in
the Marier Spheres algorithm. If elements arrive asynchronously, out of sync, as distinct and seperate
events, this recomputation could be needlessly re-triggered numerous times in a short
period. The solution requires some kind of synchronization mechanism to be employed
so that the recomputation is only triggered once. In patcher languages like Max and Pure Data,
this is typically achieved using a "pack" object, which only triggers its output when the
left-most input is received. While more sophisticated synchronisation methods could
certainly be imagined, this kind of naive scheduling by hand is typical, and
while frequently annoying and error-prone, sufficient.

Another challenge that arises is that dataset-oriented approaches introduce the need for
offline treatment of data in the form of dataset recording mechanisms, as well as ideally
tools for editing, curating, annotating, combining, and post-processing datasets.
Many dataset-oriented techniques are provided to users as add-ons
for dataflow programming environments, especially Max/MSP and Pure Data.  In
these cases, facilities for recording data are often noticeably bare.  The user
may have to write code or add objects in order to control recording, training,
and running the model, as well as sharing datasets between models, or
performing any further analysis or processing on recorded datasets
\cite{bevilacqua2005nime_mnm, gillian2011sec, francoise2014nime_xmm, kiefer2014esn,
Bullock2015}.  This code may need to be instantiated multiple times if multiple
models are used.  There may or may not be simple unified facilities for
exchanging data between models.  Furthermore, processes used to manipulate
signals should be available for the
manipulation of data points and datasets, as well as allowing dataset tools
to sit in the dataflow network.  This would further
enhance the integration between dataflow and dataset approaches, and make for a
remarkably versatile and flexible mapping design environment, but is unfortunately
not an affordance that is made readily usable in any environment surveyed.

Far from providing highly usable integration of datasets into data flow environments,
most existing tools are tools are
deliberately incomplete.  At least one tool surveyed actually requires
datasets to be recorded and models trained in MATLAB before exporting the model
in a format where it can be imported and run in Pure Data
\cite{cont2004pd-ann}.  This affords users with maximum flexibility, but
provides relatively low usability, especially for users that are not already
confident in the programming environment.  This effectively makes these
distributions completely inaccessible to potential users who lack a high level
of technical skill and confidence.

There is a significant gap in the papers surveyed for a tool that provides good
usability for non-technical users and an unified workspace providing for both
dataflow- and dataset-oriented mapping design approaches.  Particularly in
light of the unifying conceptual model proposed in
\cref{section:endpoints-and-connections}, dataset-oriented approaches may be
best utilized to enhance dataflow-oriented design tools.  In addition, dataflow-oriented
approaches could be gainfully employed to post-process datasets, providing bidirectional
enhancement to these two approaches. Future tools may
improve on existing efforts by providing readily usable tools for recording, processing, analysing,
manipulating, curating, and otherwise handling training datasets, including
processing them with the available dataflow operations, as well as allowing
dataset-oriented algorithms to exist as devices or complex connections in the dataflow graph.

Furthermore, higher-order mapping techniques have been explored in relatively
few tools to date.  Although some such approaches are relatively easy to
achieve in most dataflow design tools, others may be more involved due to
incompatibilities in the way data is represented in different approaches.
Ideally, any dataset, model, and whole mappings should have the ability to be
represented as a data point, a dataset, or a group of endpoints.  Adding such
affordances would open up the possibility to more easily explore a wide variety
of higher-order mapping approaches.

\section{Thought Experiments}

This section considers some hypothetical mapping designs to illustrate the usefulness
of an integrated design approach incorporating both dataflow- and dataset-oriented
mapping design tools, and particularly of higher-order mapping structures.
The thought experiments consist of proposed mappings
for real and imagined DMIs.

\subsection{A Higher-Order Mapping for T-Stick}

The T-Stick is a digital musical instrument invented by Joseph Malloch
\cite{Malloch2008} that consists of a PVC pipe shell on which is placed a
force-sensitive resistor (FSR) that measures squeezing pressure applied by the
player, an array of capacitive touch sensors that sense where the player's hand
contacts the pipe, and a magnetic-inertial measurement unit from which the
orientation, angular velocity, and linear acceleration of the instrument can be
estimated. In addition, one tactile push button is found on one end of the
pipe. In typical compositions for the T-Stick, a sequence of different dataflow mappings
are employed over the course of the piece, activated by a timer or push of the
button. This is convenient insofar as it allows the performer to avoid
accidentally activating unintended sounds at the wrong moment, since the
excitatory affordances of the instrument are limited to those of the current
mapping mode. This is a common approach for many DMIs, which is appropriate for
a composed piece that should generally be performed in more-or-less the same
way with similar results at each performance.

An alternative approach, more appropriate for e.g. improvisation, would be to
allow all mapping modes of the instrument to be quickly accessed from one main
mapping. The simplest way to achieve this might be to allow each mapping mode
to be accessed as a preset, e.g. by pushing the button to move to the next
preset. This implementation unfortunately doesn't provide the player with much
fluidity to move between interaction modes.

By integrating dataset-oriented tools with the various dataflow modes of the
instrument in a higher-order structure, a more flexible design can be achieved.
Suppose that the player has a few mapping modes in mind: a rain-stick mode,
where the instrument is gripped by both hands, one on either side, and a
granular noise is synthesized depending on the change in tilt and acceleration
of the instrument; a metal bar mode, where the instrument is gripped with one
hand on the end of the pipe, and sudden impulsive accelerations from jabbing or
tapping the instrument and brushing gestures on the touch sensor are used to
excite a synthesized model of a metal bar; a light-saber mode, where the
instrument is held with both hands on one end, and squeeze pressure,
orientation, and rotation influence whirring electronic synthesized sounds; and
an abstract granulator, where the instrument is held with one hand in the
middle of the pipe, and the location of a touch by the other hand selects
from a sample buffer to be granulated, while orientation and acceleration
modulate filtering and reverberant effects.

Notice that by construction, each mode can be differentiated based on the position
of the player's hands. Supposing that the output from the touch sensors for each
of these positions is recorded as a sequence of demonstration data for the
geometric model of a preset-interpolation algorthm, or similar uni-modal
learning model, then a vector of weights, one for each mapping mode, can be
computed and used to drive a higher-order mapping that integrates all four modes.

\subsection{1-Bit Whistle}

The previous example considers the T-Stick, an input device with numerous
sensors and a rich established gestural vocabulary, controlling a synthesis
engine which is presumed to be comparably rich and flexible. To motivate for
the general utility of combining dataset and dataflow tools, consider the
following example that considers an imagined instrument, the 1-Bit Whistle,
which consists of a breath pressure sensor, a handful of buttons arranged to
support a particular pitch fingering system (e.g. \cite{west2022nime_pfs}), and
a 1-Bit synthesizer \cite{troise2020thesis, troise2020_1-bit}. Even with such a deliberately constrained
design, both dataflow and dataset-oriented techniques prove useful.

Dataflow approaches naturally lend themselves to implementing models of derived
endpoints in multiple layers of processing when there is a clear or well known
way to derive interesting signals from those available. Breath pressure sensors
typically produce a non-zero voltage at ambient atmospheric pressure, which can
be reduced by inhaling through the sensor, allowing both blowing and inhaling
pressure to be measured. By high-pass filtering the breath pressure signal,
rapid changes associated with strong or sudden articulations of the airstream
can be derived. Breath pressure sensors often capture frequency data up to as
high as 10 kHz \cite{scavone2005nime_breath-pressure-frequency}; thus spectral analysis
can be employed to recognize when the player is vocalizing into the whistle,
with what amplitude, and at what pitch. A datasets-oriented tool such as a
Gaussian mixture model could be used to estimate the player's vocal formants
(caused by the posture of the oral cavity) when vocalizing into the pressure
sensor. Thus, using a combination of dataflow and datasets oriented techniques,
a variety of information can be extracted from the breath pressure signal that
isn't obviously available at first glance, including transient articulations,
growl amplitude, pitch, and format, as well as the basic pressure signals for
exhalation and inhalation.

A 1-bit synthesizer \cite{troise2020thesis, troise2020_1-bit} naturally lends
itself to producing pulse waves with pulse-width-modulation (PWM). Numerous
pitch-based effects are possible, such as vibrato, glissando and portamento,
trills, and grace notes. In addition, using pitch modulations faster than a
human can typically perform, a transient can be emphasized by very briefly
modulating the frequency at the beginning of a note e.g. starting up one octave
and dropping down after 5-20 milliseconds, and chords can be approximated by
rapidly arpeggiating their tones. Given only 1-bit, amplitude control appears
impossible. However, thinner pulse-widths also result in a lower perceived
amplitude. This can be exploited to provide some amplitude control. Automating
this effect, a feedback-delay-like echo effect can be simulated by repeating
note onsets at lower pulse-widths. Thus, with an appropriate dataflow network
implementing a higher-level semantic model, even a characteristically 1-bit
synthesizer could be considered as having controls for pitch, transients, super
fast arpeggiation, echo, and pulse width modulation.

Breath pressure could be mapped to pulse width to provide a basic amplitude
control, with sudden variations in pressure (as detected by a high-pass filter)
causing pitch variations, including octave-up transients on strong
articulations, or causing vibrato with slower pressure variations. Convergent
mappings may also prove useful. For instance, the current octave could be used
to adjust the base mappings just described, so that each octave has a unique
characteristic pulse-width when the breath pressure exceeds a base threshold.
Alternatively, the breath pressure to pulse width mapping could be implemented
using a breakpoint envelope, so that the control points of the envelope could
be adjusted continuously depending on the octave. The transient-emphasizing
modulation could be similarly adjusted to give each octave its own character,
e.g. by adjusting the times in a modulation table.

There are a few possibilities for the vocalization signals that spring to mind. On
a basic level, the presence or absence of vocalization could be used to enable
a super-fast arpeggio, with the speed and/or number of notes depending on the
amplitude, pitch, and/or formant of the vocalization. Alternatively, depending
on the formant, various different effects could be activated, such as portamento
or pulse-width-based echo.

\subsection{Thought Experiments Summary}

The mappings described are purely hypothetical, and at a necessarily high level.
Actual implementation of these designs would almost certainly reveal issues
that are obvious when interacting with the system but imperceptible at the
level of this conceptual discussion. Nevertheless, the examples illustrate
some of the advantages that arise from integrating dataflow- and dataset-oriented
approaches to mapping design. In particular, recognizing where dataflow mappings
can be viewed as being determined by datasets, such as in breakpoint envelopes
and modulation matrices, and then modulating those datasets in higher-order
mapping structures, is a highly versatile approach that is in general not well
supported by existing dataflow environments without extensive manual implementation.

\section{An Imaginary Mapping Tool}

Let us imagine a deeply integrated mapping authoring tool,
called the Workbench, drawing on the framework and tools described thus far. Such
a tool should facilitate seamlessly working with dataflow and dataset oriented approaches
together, including the design of higher-order mapping structures that blur the lines
between these approaches.

The Workbench supports two complementary workflows: authoring mapping networks,
and recording and editing datasets. Under the umbrella of these two modes of interaction,
numerous different levels of abstraction and alternative perspectives are supported.
The integration of these two modes of interaction runs sufficiently deep so that
operations in one mode seamlessly blend into equivalent actions in the other\footnote{a
likely application for bidirectional transformations..?!}.

\subsection{Throughpoints}

Unfortunately, while the conceptual model of endpoints and connections seems
adequate for thinking about the semantic encoding of a mapping design, it is
somewhat at odds with the practicalities of an implementation. Dataflow is the
overriding conceptual model of implementation, where endpoints are sources or
destinations for signals. Dataset-trained models are naturally treated as
devices with endpoints, but this is at odds with the realization that the true
endpoints of a trained model are implicitly encoded in the dataset it learns
from. A simple extension of the model is to introduce a distinct kind of
endpoint that reflects the semantically neutral character of the inputs and
outputs of a trained model. Unlike the conceptual endpoints desribed above,
which have particular meaning and practical interpretation in the context of
the network, the inputs and outputs of a trained model at runtime are mere
reflections of the incoming signals and influenced parameters connected to the
model. These inputs and outputs are further distinguished from typical
endpoints in that they are agnostic about the shape and structure of the
signals flowing through them. As long as the sources and destinations conform
to the specifications given by the dataset (or vica versa, as long as the
dataset conforms to the specification given by the input and output signals
connected through the model), dataset-trained models can generally adapt to
arbitrary input-output dimensionality ($N$ to $M$). Let us term
dataset-oriented devices' inputs and outputs, and other inputs and outputs
having these qualities, as \emph{throughpoints}, since data flows through them
without a shift in its semantics, unlike endpoints which imply a change in the
interpretation of the data flowing through them. So that we can refer to the
set containing both throughpoints and endpoints, let us also introduce the term
\emph{ports} for this set.

Dataset-oriented tools are not unique in possessing throughpoints. Another
example of a throughpoint might be a binding to a network communication
protocol such as libmapper. When the endpoints of a device are connected to the
libmapper network, this connection should ideally adapt to the endpoints of the
device and reflect them unchanged onto the network. Other kinds of binding glue
are similar, such as bindings to instantiate a device as a Pure Data external
or DAW plugin.

\subsection{Authoring Mapping Networks}

In the mapping authoring mode, the mapping network is visualized in a more-or-less typical
boxes-and-lines diagram. Endpoints are visualized with small circular icons. Connections
are drawn as lines from one endpoint to another. Boxes are drawn intersecting endpoints
and grouping together connections. Boxes are used to represent devices, in the sense of
dataflow devices. Devices can have arbitrary and heterogeneous number of inputs
and outputs.

In addition to endpoints, devices may also possess throughpoints, visualized
as rectangular boxes on the edges of devices. Whereas endpoints may only
be connected to endpoints with compatible type (e.g. dimensionality), and only one
signal may be connected to one endpoint, any number of arbitrary signals can be connected
to a destination throughpoint, and it is expected to automatically adapt to accept them.
Similarly, a source throughpoint may be connected to an arbitrary number of destinations
of arbitrary type, and the device is expected to automatically adapt its behavior
to support these connections appropriately. Throughpoints are mainly associated
with devices that implement dataset-oriented approaches, as well as bindings to external
protocols.

Devices themselves can be implemented using the mapping authoring language. Such devices are
called subnets, as they contain a part of the overall mapping network. They function as subpatches
or abstractions, similar to in patcher languages like Max and Pure Data. The viewport
can be dynamically zoomed into and out of such devices, or they can be expanded, allowing
their internal implementation to be viewed simultaneously with the higher-level patch.
In addition, the highest-level view can always be zoomed out of and treated as a subnet,
allowing a higher level of abstraction to be introduced at any time, and facilitating the
development of higher-order mappings.

At the lowest level, primitive processing elements such as arithmetic are
represented by devices that cannot be examined in this way. Perhaps these
primitive devices could be visualized as black boxes, in contrast to other
devices seen as transparent boxes. External algorithms, e.g. implemented in
C/C++, could be shown in the same way. Some black boxes are notable because
they autonomously produce source data or provide final destination ports that
influence external processes such as producing sound through the
analog-to-digital converter, or visualizing video signals on a display.
Original source ports in particular have special significance when recording
datasets, described later.

Scalar signals that are input-to or output-from a device at the same moment in time
can be viewed as vector-valued signals.
The display can be optionally set to visualize them as such, grouping together these
otherwise separate connections into thicker cables and simplifying the visual impact and
interpretation of the display and bringing it more in line with the interpretation
favored in a datasets-oriented approach. Similarly, vectors with the same dimension
that simultaneously arrive or depart from a box can be viewed as columns or rows of
a 2D matrix-valued signal. As will be elaborated momentarily, higher-dimensional matrix structures must also
be admitted to facilitate higher order mappings,
and also allowing for e.g. video processing networks to be represented.
Primitive devices are provided to explicitly annotate grouping, ungrouping, and general recombination of
multi-dimensional signals, and defining the synchronization semantics of such
combinators.

An inspector pane is provided allowing a connection, device, or port to be examined.
This could show the metadata of the object, allow its variables to be edited,
and visualize the signals flowing through, into, or out of the object.
The variables of a connection or port are simply the signals flowing through them.
An unconnected port can have its value set as a constant.
Devices similarly present their ports' values as variables in the inspector.
In addition, any unconnected ports inside the subnet implementing the device
can also be viewed and their constant value modified. Such unconnected inner ports
can also be promoted to being ports of the subnet, allowing signals and destinations of the parent
subnet to be connected to them.

Connections are inherently passive; they are not allowed to perform any computation,
and merely represent data flowing from one device to another. As such, basic
mapping operations such as affine transformation must be implemented using primitive
devices, or more usefully using standard devices provided by the Workbench for the most common use-cases.

For example, when a simple connection is made, by default it could create an instance of a modulation
matrix subnet, allowing basic affine transformations to be quickly implemented,
and providing useful representation of such basic mappings in a way that is amenable
to the higher-order extensions to the basic modulation matrix. When additional connections
are made, if they can be represented as part of an existing modulation matrix, they are
automatically grouped together with it.

\subsection{Representation of Datasets in the Network}

There are four kinds of datasets, in line with the taxonomy given by Francoise 
et al \cite{francoise2014nime_xmm}: snapshots (instantaneous unimodal), 
sequences (temporal unimodal), demonstrations (instantaneous multimodal), and 
automations (temporal multimodel). Automations can be decomposed into two 
sequences, one made from all the source signals, and one made from all the 
destination parameters. Similarly, demonstrations can be decomposed into two 
snapshots. Sequences can further be decomposed into snapshots (one per time 
sample). Inversly, snapshots or sequences can be combined cross-modally to make 
demonstrations and automations, respectively, and snapshots or demonstrations 
can be combined temporally to make sequences or automations, respectively.

Snapshots, sequences, demonstrations, and automations (datasets, generically) 
can be viewed as destination endpoints of dataset-oriented subnets, and they 
can optionally be exposed as such on the edges of the subnet to facilitate 
higher-order mappings. Such exposed datasets can be decomposed at any time, so 
that e.g. an exposed sequence can be treated as a collection of exposed 
snapshots, or vica versa.

Consequently, datasets themselves may flow as data through the mapping network. 
As such, snapshots and demonstrations may be viewed as vectors, and sequences 
and automations may be viewed as matrices. More generally, sequences and 
automations could be viewed as higher-dimensional snapshots and demonstrations 
in higher-order structures, and sequences and automations could also be 
composed of a time-sequence of matrices rather than vectors. Various detailed 
facilities would be required to help the user visualize and conceptualize data 
at these varying levels of structure and abstraction, building on the basic 
operations of zooming into and out of subnets, and viewing groups of scalars as 
vectors, groups of vectors as matrices, and vica versa, and so on.

\subsection{Authoring Datasets}

In the dataset authoring mode, subnets whose data the user is interested in
recording and editing (marked as such from the mapping authoring mode) can be
viewed as tracks in a digital-audio-workstation-informed editor. Subnets which
require a dataset in order to function are automatically added as tracks to the
recorder/editor. Each track records a subset of the signals in the network,
with various selections possible. In the most common case, signals input to a
subnet, and destinations influenced by it, are selected for recording. Other
possibilities include recording the free endpoints of a subnet (both input and
outputs, or one or the other), its incoming signals or influenced destinations
only, or some other arbitrary combination of signals deemed to be of interest.

Recordings can be freely sliced, spliced, cropped, duplicated, etc., both
temporally, and across their multiple signal dimensions. Internally, all such
edits are represented by automatically generated mapping subnets that processes
the raw recorded data. The raw recorded data, in turn, solely consists of
original source signals and final destination parameters associated with black
boxes such as external devices. This is recorded with the temporal
characteristics in which it originally arrived. All synchronization, grouping
into vectors, etc., is encoded in a snapshot of the mapping network state at
the moment of the recording. This mapping network can be arbitrarily (and
non-destructively) edited by the user to modify the dataset, or it can be kept
in sync with edits made to the mapping network, so that a recording needn't be
discarded due to a later change to the mapping network. As long as the
characteristics of the external devices remain consistent, recordings remain
valid, no matter what edits are made. This also allows the signals of interest
in a dataset to be arbitrarily modified at a later time, e.g. if the signals
input to a preset interpolator are changed at some point in the design, the
training data need not necessarily be regenerated.

This also allows recordings to serve as useful prototyping and collaboration
tools. For example, a recording could be taken when in person with a
collaborator who has a particular input device not usually accessible to the
mapping designer. Later, the mapping designer could replay this recording, e.g.
using a built in playback subnet, and continue to edit the mapping network
without the collaborator's presence or participation.

Finally, recordings need to be associated with the devices that use them. This
can be done from the multi-track editor by enabling or disabling the recordings
on a track; the enabled recordings are automatically fed to the subnet that
consumes them as constant values for its dataset endpoint. This is, again,
internally represented by an automatically generated subnet that can be
examined and edited by the user, eventually facilitating the development of
higher-order mappings.

\subsection{Summary}

The Workbench is based on a unified underlying dataflow model, a graphical
programming language where primitive black box processors are combined in
dataflow graphs to create high level and high order programs that implement
mappings. Signals are allowed to have arbitrary dimensionality and structure,
so that datasets can be represented as signals in higher-order mapping
structures. High-dimensional signals can be readily decomposed into
lower-dimensional subcomponents, and lower-dimensional signals can be combined
into higher-dimensional signals. In the case of datasets, time itself is just
another dimension of the dataset-viewed-as-signal, and the dataset-signal can
be decomposed along the time dimension or any other dimension. Such operations,
and all other edits the user may want to make to a dataset, are
non-destructively represented using the unified dataflow language. All recorded
datasets are derived from original raw source signals, so that as long as the
external signal sources remain consistent, all recordings can be repurposed and
reinterpreted anywhere in the network. When datasets are utilized on the
network, this too is represented by the underlying dataflow language, with the
dataset used as a constant offset value for a multi-dimensional signal. The use
of throughpoints allows a natural representation of the dimensionally-agnostic
ports of dataset-oriented mapping design tools, as well as bindings to external
environments.

The hypothetical key features of the Workbench, besides being highly usable,
attractive, efficient, performant, and so on, would be its representation of
dataset-oriented approaches, and of dataset editions, within a pure dataflow
paradigm, by using throughpoints and primitive devices that facilitate working
with flows of data having arbitrary dimensionality and structure.
These are the underlying qualities that would enable a tight integration
of dataset-oriented approaches with dataflow-oriented techniques, as well as
opening up many possibilities for higher-order mappings.

As a final provocation, perhaps the edition of the network and its datasets
could itself be represented by a meta-programming subnet which accepts user
actions as its inputs and produces new states of the network and its datasets
as outputs. A recording of the users actions would then constitute a complete
history of the mapping design, and could be used to generate a version history,
allow non-linear undo, forking of the time-line, and so on.

\section{Executive Summary}

\begin{displayquote}
Several mapping tools have been proposed to facilitate the prototyping of DMIs.
These mostly fit into two main categories: datasets (interpolations) and data
flow (parameter mapping). Each has its own strengths and weaknesses, their
choice likely impacting the final mappings done.
\end{displayquote}

I conducted a review of mapping authoring tools, and described them in terms of
the main objects of interest as being oriented toward dataflow or datasets. I
reflected on how these different tools could be combined using higher-order
mappings. In considering the design of mappings with dataflow and datasets and
the ways in which semantic associations are recorded in these mappings, and
drawing on prior characterizations of mappings as having multiple layers, I
proposed to distinguish between cross-domain or cross-modal mappings, which are
mappings that associate from semantically meaningful endpoints to other
independently semantically meaningful endpoints, and models, which are
processes that derive new semantically meaningful endpoints from existing ones.

\begin{displayquote}
Describe in detail 2 examples of tools in each category and elaborate on the
possible impacts of each tool in the mappings obtained with them
\end{displayquote}

I described in detail the modulation matrix, and other common techniques used
in dataflow mappings. I described in detail inverse weighted distance and
Marier spheres preset interpolation algorithms, as examples of dataset-oriented
mapping design techniques. I reflected on how these tools influence the design
process and the kinds of mappings that can be readily produced with them.
Furthermore, I provided a broad overview of numerous other tools presented in
the literature, and contemplated their orientation towards data and their
possible combinations, broadly.

\begin{displayquote}
Discuss how a hypothetical solution could combine both approaches and present a
pseudo-code that could implement such a tool for one specific NIME context
(programming environment, interface, synthesis method). Comment on whether
these types of hybrid tools could be generalized to multiple music interaction
contexts.
\end{displayquote}

I described in detail two designs which could leverage a combination of both
dataflow and datasets in their mappings, based on two deliberately very
different interaction contexts involving different kinds of sensors, different
number of degrees of freedom, and contrasting approaches to sound synthesis. In
addition, I described a hypothetical integrated DMI development environment
which would allow the seamless interoperation of dataflow and datasets by
providing a unified underlying conceptual model that flexibly hybridises these
dissimilar approaches. The main features of this conceptual model are the
notion of "throughpoints", and the treatment of arbitrary datasets as just
another signal that can flow through the processing network. I introduced the
term "throughpoints" to characterise ports of dataset-oriented algorithms that
adapt fluidly to the dimension and characteristics of their inputs, based on
training data, in contrast with "endpoints" that have their own independent
semantic associations. The treatment of datasets as data, combined with
throughpoints, allows dataflow techniques to be applied to datasets, unifying
the two approaches.

\end{document}
